\section{Arbitrary and repeated row arguments}\label{sec:repeated-moments}

To obtain repeated rows at arbitrary arguments, add a fused pair of sites for each desired argument and remove the pairs after applying the distinct-site formula. The following collision calculation justifies that removal for the scalar trace.

\begin{theorem}[All row moments between the vacuum caps]\label{thm:row-moments}
Let $N,h\ge0$ be integers. With ordered site variables $\mathbf b=(b_1,\ldots,b_N)$, the right and left vacuum caps $\Psi,\Lambda$, the scalar trace $\tau$, and the Laurent polynomial part $P$ of Sections~\ref{sec:vacuum}--\ref{sec:site-moments}, the following is an identity of meromorphic functions of the site variables, nonzero twist $K$, and row variables $u_1,\ldots,u_h$:
\begin{equation}\label{eq:all-row-moments}
 \Lambda\prod_{a=1}^h T^{K^{-1}}(u_a;\mathbf b)\Psi
 =\frac{\tau\!\left(\prod_{a=1}^h P(u_a)/d_{\rm full}(u_a)\right)}{\tau(1)}.
\end{equation}
At a common-twist collision, the trace is the holomorphic specialization near a generic diagonal supplied by Proposition~\ref{prop:common-twist}; colliding roots are retained by that continuation. Pointwise evaluation is at regular values of the meromorphic expressions, including repeated row arguments. For $N=0$ the empty site and root conventions apply, and for $h=0$ both sides are one.
\end{theorem}

We prove the theorem by a two-site reduction. The matrix side reduces by fusion. On the scalar side, we show that both traces have at most a simple pole, with residues equal to a common nonzero factor times the corresponding smaller traces. The denominator residue is nonzero generically, so that factor cancels in the normalized ratio.

\subsection{A scalar two-site reduction}

Let $\mathbf b_{\rm red}$ be the ordered list of $n\ge0$ sites obtained by deleting two neighboring sites $x,y$, in that order, from a larger list $\mathbf b_\epsilon$, where $x=y+3+\epsilon$. All other sites are generic. Use subscripts $\epsilon$ and ${\rm red}$ on $\tau,P,d_{\rm full},W,D_i,J$ to distinguish the two lists; a subscript $0$ on a regular coefficient denotes its specialization at $\epsilon=0$.

\begin{lemma}[Two-site reduction of the normalized trace]\label{lem:scalar-two-site-fusion}
Fix a generic common twist for both lists. For any finite list of row arguments $u_a(\epsilon)$ with limits $u_a(0)$, assume that for every $a$ both limiting denominator evaluations are regular and nonzero:
\[
 d_{{\rm full},0}(u_a(0))\ne0,
 \qquad d_{{\rm full},{\rm red}}(u_a(0))\ne0.
\]
After shrinking the $\epsilon$-disk, the nearby evaluations used below are then regular and nonzero. Put
\[
 f_\epsilon=\prod_a\frac{P_\epsilon(u_a(\epsilon))}
                         {d_{{\rm full},\epsilon}(u_a(\epsilon))},
 \qquad
 f_{\rm red}=\prod_a\frac{P_{\rm red}(u_a(0))}
                         {d_{{\rm full},{\rm red}}(u_a(0))}.
\]
Then, at generic values of the remaining parameters,
\begin{equation}\label{eq:scalar-two-site-fusion}
 \lim_{\epsilon\to0}\frac{\tau_\epsilon(f_\epsilon)}{\tau_\epsilon(1)}
 =\frac{\tau_{\rm red}(f_{\rm red})}{\tau_{\rm red}(1)}.
\end{equation}
The identity extends meromorphically in the remaining parameters. Empty products are allowed. The row arguments may vary with $\epsilon$ subject to the stated regularity.
\end{lemma}

\begin{proof}
We first bound the order at the merging sites and identify the contributing root cluster. We then compute its residue and restrict the polynomial insertions.

\paragraph{Preparing the site collision.}
Keep the \(n+2\) twists \(K_i\) independent during the transverse calculation, while the polynomial insertion still uses the single variable \(K\). We calculate residues coefficientwise for a Laurent-polynomial insertion \(F_\epsilon\) in the roots whose coefficients are rational in the exponential parameter coordinates and jointly holomorphic in the site variables (including \(\epsilon\)), the twists \(K,K_i\), and the row variables near the chosen collision point and generic diagonal, after inverting the stated nonzero scalar denominators. Both \(F_\epsilon=1\) and \(F_\epsilon=f_\epsilon\) have this form. The calculation will show that \(\epsilon\tau_\epsilon(F_\epsilon)\) has a regular limit and will express that limit as a sum of smaller independent-twist traces. We keep the row variables independent in this calculation and substitute their prescribed limits at the end.

The multiplicative site coordinates remain in a fixed compact annulus and have distinct shift orbits for all sufficiently small \(\epsilon\ne0\). We may therefore use the uniform punctured-neighborhood conclusion in the proof of Proposition~\ref{prop:common-twist}: one neighborhood \(U\) of a generic diagonal, avoiding the subset-product toric-end hypersurfaces for the large trace and for every smaller trace obtained by deleting two twist coordinates, works throughout this punctured site family. With the remaining site and row parameters near their chosen regular values, \(\tau_\epsilon(F_\epsilon)\) is holomorphic there, so no other polar divisor can approach the collision inside \(U\). Once at-most-simple order is proved for generic independent twists, \(\epsilon\tau_\epsilon(F_\epsilon)\) is therefore regular jointly near the generic diagonal: one may use holomorphic extension off codimension two, or cancellation of every possible denominator divisor in the rational function. This will carry the independent-twist residue identity to the diagonal. The transverse calculation itself is made off all other exceptional sets of twists on \(\epsilon=0\).


\paragraph{The only contributing cluster.}
Use the compactification argument of Lemma~\ref{lem:residue} over a small disk of site values, counting moving \(x\) as one more configuration coordinate relative to fixed anchors; end limits remain excluded. Compute order of \(W\prod dt_i\wedge d\epsilon\) relative to log poles. Single-site and unanchored depths give the earlier bounds for bare \(W\) (a cluster drifting with the moving site but not \(y\) counts as anchored), whereas until separation of \(x,y\) the cluster containing both gives
\[
 B-m_{-1}-m_1-m_2-m_4+1
\]
with offsets relative to $y$, since the log discrepancy is now one larger. The following claim gives both the order bound and its possible equality case.

\begin{claim}[Order at the merging sites]\label{claim:merging-site-order}
For the offset counts of a cluster containing both sites, put
\[
 B=\frac12\sum_{k\bmod8}(m_k-m_{k+1})^2
       +\sum_{k\bmod8}m_km_{k+2}.
\]
Then
\[
 B-m_{-1}-m_1-m_2-m_4+1\ge0,
\]
and equality requires unit support $\{1,2\}$.
\end{claim}

\begin{proof}
If the minimum count is positive, the distance-two products already suffice; one can use the terms starting at each of the four anchored indices. Suppose henceforth that at least one count is zero, and write
\[
 a=m_{-1},\quad b=m_1,\quad c=m_2,\quad f=m_4,
 \qquad M=\max_k m_k.
\]

\emph{A doubly occupied anchored pair.}
If either pair $(a,b)$ or $(c,f)$ has both counts positive, use $B\ge M+ab+cf$. When $a,b>0$, the inequality $(a-1)(b-1)\ge0$ gives $a+b\le ab+1$, and analogously $c+f\le cf+1$ when $c,f>0$. A singly occupied pair contributes at most $M$. If both pairs are doubly occupied, use also $2\le M+1$. In every case this gives
\[
 a+b+c+f\le M+ab+cf+1\le B+1;
\]
the first inequality is strict when the other pair is empty. Equality thus requires the other pair to be occupied too, equality in the variation bound, and no distance-two products beyond $ab,cf$. In particular the support must be contiguous. This is impossible: contiguous support containing the two indices of the first doubly positive pair and an index of the other either gives an additional unwanted product, or makes both pairs doubly occupied with unwanted products.

\emph{Neither anchored pair is doubly occupied.}
Use $B\ge2M-1$; the empty case is direct. The proof of that bound in Proposition~\ref{prop:common-twist} gives strictness for $M\ge3$. At $M=3$, the four gradients give an integral bound that could be sharp only if both first neighbors were positive, adding their product. If a second neighbor is positive, the variation and its product give at least $2M$ altogether.

For $M=2$, equality only needs to be excluded when two peaks have height $2$. All their second neighbors must vanish. The variation-square mass on the four edges around one peak is already at least $2$, leaving at most $1$; hence no other height-$2$ peak can lie outside its first neighbors, since the remaining path starts and ends at zero. If one first neighbor equals $2$, the variation there is already at least $3$ and the other first neighbor must be unit, adding their product. Thus the two distinct anchored pairs cannot both supply height $2$.

For $M=1$, the unit equality sets for $B=M$ give precisely the asserted pair $\{1,2\}$. This proves the claim.
\end{proof}

Over generic independent twists all levels are transverse on boundary strata and solutions over \(\epsilon=0\) lie on only one boundary component (codimension at most one). A nonzero log order cannot contribute to the limiting trace of \(\epsilon W f\) (residues divided by log Jacobian understood): locally use level values as coordinates with one further coordinate \(v\), \(\epsilon=v^m\) after absorbing a unit. Neutrality and log order zero are thus necessary. Clusters off the merging orbit are strictly excluded by the argument of common-twist regularity for \(W\). In the merging orbit equality requires exactly two distinct labels \(A,B\) at offsets 1,2 at every depth containing both sites. In the local boundary coordinate this implies
\[
 \operatorname{ord}(t_A-y-1),\ \operatorname{ord}(t_B-y-2)
 \ge \operatorname{ord}\epsilon,
\]
so \(t_A=y+1+\epsilon s,t_B=y+2+\epsilon u\) with bounded \(s,u\). The scaled equations will decide whether any further singular collision occurs in these coordinates.


\paragraph{The two scaled root equations.}
Since \(S(z)S(z-3)R(z-1)R(z-2)=1\), each surviving equation tends to its equation in the regular smaller system, independently of \(s,u\). After collecting their regular nonzero factors, the two other equations have the cross-multiplied limits
\[
 C_A(s-u)=-K_A s,\qquad C_B(u-1)=-K_B(s-u)
\]
with nonzero constants depending on the surviving regular solution and the sites, but not on \(s,u\). These equations exclude the further singular collisions \(s=0\), \(u=1\), and \(s=u\): the first would force \(u=0\) and contradict the second equation; the second would force \(s=1\) and contradict the first; the third would force \(s=u=0\) and again contradict the second. For generic independent twists the equations have a unique finite simple solution. Together with a regular smaller solution it gives an actual branch by the implicit function theorem. At that solution the limiting level maps can be written
\[
 D_A^{(0)}=C_A\frac{s-u}{s},\qquad
 D_B^{(0)}=C_B\frac{u-1}{s-u}.
\]

\paragraph{The common scalar residue.}
We now compute the factor contributed by one ordered choice of the deleted root labels \((A,B)\). In the ordinary angle coordinate \(\epsilon\), substitution into the scalar functions \(V,G\) of Lemma~\ref{lem:residue} gives
\[
 \lim_{\epsilon\to0}\epsilon^3s(u-1)(s-u)
 V(t_A-x)V(t_A-y)V(t_B-x)V(t_B-y)G(t_A-t_B)
 =\frac{[l]}{8[2l]^3[3l]}.
\]
Each remaining site \(b\) supplies \(V(y+1-b)V(y+2-b)\), whereas every surviving root factor cancels by
\(G(z-1)G(z-2)V(z)V(z-3)=1\). Define
\[
 \kappa(\mathbf b_{\rm red};y)
 :=\frac{[l]}{8[2l]^3[3l]}
   \prod_{b\in\mathbf b_{\rm red}}V(y+1-b)V(y+2-b).
\]
This is a nonzero meromorphic function at generic sites. It may depend on the collision center \(y\) and every remaining site, but is independent of the surviving roots, all twists, and the ordered choice of deleted labels. The full weight consequently satisfies
\[
 \epsilon^3W_\epsilon\longrightarrow
 \frac{\kappa(\mathbf b_{\rm red};y)W_{\rm red}}
      {s(u-1)(s-u)}.
\]

The scaled Jacobian has exactly the same factor in \(s,u\). Indeed
\[
 \partial_{(s,u)}\begin{pmatrix}\log D_A^{(0)}\\ \log D_B^{(0)}\end{pmatrix}
 =\begin{pmatrix}
  (s-u)^{-1}-s^{-1} & -(s-u)^{-1}\\
  -(s-u)^{-1} & (u-1)^{-1}+(s-u)^{-1}
 \end{pmatrix},
 \qquad
 \det=\frac1{s(u-1)(s-u)}.
\]
The surviving level maps are independent of \(s,u\) in the limit, so the full derivative matrix in the scaled coordinates is block triangular, even though \(C_A,C_B\) may depend on the surviving roots. Since \(t_A,t_B\) each contribute a factor \(\epsilon\) to the coordinate change, and simultaneous row and column reorderings introduce no sign,
\[
 \epsilon^2J_\epsilon\longrightarrow
 \frac{J_{\rm red}}{s(u-1)(s-u)},
 \qquad
 \epsilon\frac{W_\epsilon}{J_\epsilon}\longrightarrow
 \kappa(\mathbf b_{\rm red};y)\frac{W_{\rm red}}{J_{\rm red}}.
\]

For a coefficientwise limit \(F_0\), write
\[
 \iota_{A,B}^*F_0:=
 \left.F_0\right|_{t_A=y+1,\ t_B=y+2}
\]
for its restriction to the contributing root cluster. Write \(\tau_{\rm red}^{\widehat A,\widehat B}\) for the smaller trace with the surviving independent twists. Summing the branch limits gives, for generic independent twists,
\[
 \operatorname{Res}_{\epsilon=0}\tau_\epsilon(F_\epsilon)
 :=\lim_{\epsilon\to0}\epsilon\tau_\epsilon(F_\epsilon)
 =\kappa(\mathbf b_{\rm red};y)
  \sum_{A\ne B}\tau_{\rm red}^{\widehat A,\widehat B}
                   (\iota_{A,B}^*F_0).
\]
The sum is over the \((n+2)(n+1)\) ordered choices of \((A,B)\). The simple-order and joint-regularity argument above continues this rational identity to a generic diagonal as an equality of these residue limits. For \(F_\epsilon=1\) or \(F_\epsilon=f_\epsilon\), the restrictions are symmetric in the surviving roots, so all summands become the same smaller trace after relabeling. Thus the common nonzero trace residue factor is
\[
 \mathcal C_{\rm del}(\mathbf b_{\rm red};y)
 :=(n+2)(n+1)\kappa(\mathbf b_{\rm red};y),
\]
and at a generic common twist
\begin{equation}\label{eq:scalar-fusion-residue}
 \lim_{\epsilon\to0}\epsilon\tau_\epsilon(F_\epsilon)
 =\mathcal C_{\rm del}(\mathbf b_{\rm red};y)\,
   \tau_{\rm red}(\iota^*F_0),
 \qquad F_\epsilon=1\ \text{or}\ f_\epsilon.
\end{equation}
Here \(\iota^*\) denotes any ordered choice after relabeling the surviving roots. In particular, \eqref{eq:trace-normalization} gives the nonzero generic denominator residue
\(\mathcal C_{\rm del}(\mathbf b_{\rm red};y)\tau_{\rm red}(1)\).


\paragraph{Restricting the polynomial insertions.}
It remains to identify \(\tau_{\rm red}(\iota^*f_0)\) in \eqref{eq:scalar-fusion-residue}. At the collision put
\[
 d_{\rm pair}(w):=\frac{d_{{\rm full},0}(w)}{d_{{\rm full},{\rm red}}(w)}
 =[w-y][w-y-1][w-y+3][w-y+2],
 \qquad Q_{\rm pair}(w):=[w-y-2][w-y-3].
\]
Then \(\iota^*Q_0=Q_{\rm red}Q_{\rm pair}\). In every summand of \(L/d_{\rm full}\), the added \(Q_{\rm pair}\) factors cancel the two added sites using \([z+8]=-[z]\). Equivalently, as rational functions of the spectral variable,
\[
 \iota^*L_0(w)=d_{\rm pair}(w)L_{\rm red}(w).
\]

For fixed generic remaining sites and common twist, let
\(\mathcal A_{\rm red}=\mathbb C[p_i^{\pm1}:i\ne A,B]\), where \(p_i=e^{2it_i}\). Choose root and spectral monomials so that \(C_{\rm red}(z)\in\mathcal A_{\rm red}[z]\) is the polynomial form of \(Q_{\rm red}(w)Q_{\rm red}(w-1)\), with unit constant and leading coefficients. Set
\[
 \mathcal R_{\rm red}(z):=C_{\rm red}(z)(L_{\rm red}(w)-P_{\rm red}(w))
       =\sum_j r_jz^j\in\mathcal A_{\rm red}[z^{\pm1}],
 \qquad
 I_{\rm red}:=(r_j:j\in\mathbb Z)
       \subset\mathcal A_{\rm red}.
\]
Multiplying \(\mathcal R_{\rm red}\) by a spectral monomial makes it a cleared polynomial of the kind used in Proposition~\ref{prop:common-twist} and does not change its coefficient ideal. Let \(\operatorname{pp}_z\) denote the Laurent polynomial part used to define \(P\). In the formal divisions defining \(P_\epsilon\), the leading and trailing coefficients of the root denominators are monomial units on the torus, so division and hence taking \(\operatorname{pp}_z\) commute with \(\iota^*\). Since \(d_{\rm pair}P_{\rm red}\) is already a Laurent polynomial, the exact restriction needed here is
\begin{equation}\label{eq:scalar-fusion-polynomial-ideal}
 \iota^*P_0(w)-d_{\rm pair}(w)P_{\rm red}(w)
 =\operatorname{pp}_z\!\left(d_{\rm pair}(w)
                    \frac{\mathcal R_{\rm red}(z)}{C_{\rm red}(z)}\right)
 \in I_{\rm red}\,\mathcal A_{\rm red}[z^{\pm1}].
\end{equation}
Indeed formal division at zero and infinity expresses every coefficient of this finite polynomial part as an \(\mathcal A_{\rm red}\)-linear combination of the coefficients of \(\mathcal R_{\rm red}\). The divisions use the coefficients of \(C_{\rm red}\) and inverses only of its constant and leading root-monomial units.

At each allowed row value, divide \eqref{eq:scalar-fusion-polynomial-ideal} by the nonzero scalar
\(d_{\rm pair}(u_a(0))d_{{\rm full},{\rm red}}(u_a(0))\). Each normalized insertion difference then belongs to \(I_{\rm red}\), after adjoining the row-parameter scalars. Telescoping a product of such factors, including the empty product, gives
\begin{equation}\label{eq:scalar-fusion-product-ideal}
 \iota^*f_0-f_{\rm red}\in I_{\rm red}.
\end{equation}
Proposition~\ref{prop:common-twist} says that \(\tau_{\rm red}\) annihilates every Laurent-polynomial multiple of every generator of this ideal. It therefore annihilates the product difference in \eqref{eq:scalar-fusion-product-ideal}, including its contributions from all collision jets. Substitution in \eqref{eq:scalar-fusion-residue} gives numerator residue \(\mathcal C_{\rm del}\tau_{\rm red}(f_{\rm red})\) and denominator residue \(\mathcal C_{\rm del}\tau_{\rm red}(1)\ne0\). Their common factor cancels. The assumed regularity lets the row variables approach any prescribed \(u_a(0)\), proving \eqref{eq:scalar-two-site-fusion}; the stated meromorphic continuation follows from the rational identities used throughout.
\end{proof}

\subsection{Removing the auxiliary pairs}

\begin{proof}[Proof of Theorem~\ref{thm:row-moments}]
First take the desired row arguments distinct and generic. Add one neighboring pair for each argument, with its first site $x$ equal to that argument. Apply the distinct-site identity \eqref{eq:site-moment} to the enlarged list, using the subset of all these first sites, and use Proposition~\ref{prop:common-twist} to replace $L$ by $P$ in the scalar trace.

Remove the pairs successively at $x=y+3$. The denominator product at the first site satisfies $d_{\rm full}(x)\ne0$ on fusion. On the scalar side, Lemma~\ref{lem:scalar-two-site-fusion} gives the smaller normalized trace at each removal. On the transfer side, the vacuum injections and the $r=3$ transfer intertwining in Section~\ref{sec:local}, together with the vacuum fusion and left-cap contraction of Sections~\ref{sec:vacuum} and~\ref{sec:site-moments}, give the identical reduction. Thus the two sides pass through the same sequence:
\[
 \begin{aligned}
 \text{enlarged distinct-site moment}
 &\ \longrightarrow\ \text{successive two-site reductions}\\
 &\ \longrightarrow\ \text{requested row moment on }\mathbf b.
 \end{aligned}
\]
This proves \eqref{eq:all-row-moments} for generic distinct arguments. Meromorphic continuation proves it for arbitrary row arguments at regular values, including repetitions. For the empty product, the identity is the normalization $\Lambda\Psi=1$ and $\tau(1)/\tau(1)=1$.
\end{proof}
